\subsection{Baseline Results}
\label{sec:baseline}

\paragraph{Reference run.}
We evaluate the baseline method defined in \cref{sec:method:baseline}
under the shared training and evaluation protocol of
\cref{sec:exp:training,sec:exp:eval}. This run serves as the common
reference point for all later comparisons.

\paragraph{Headline numbers.}
The baseline reaches its best validation accuracy of \textbf{80.42\%}
at epoch \textbf{74}, with a validation loss of \textbf{0.9317}. Using
the checkpoint from this epoch, the model achieves a test accuracy of
\textbf{79.65\%}. These results are used as the reference values in the
later method comparisons.

\begin{figure}[!htbp]
  \centering
  % \includegraphics[width=0.85\linewidth]{baseline_curves.png}
  \fbox{\rule{0pt}{4cm}\rule{0.8\linewidth}{0pt}}
  \caption{Baseline ViT-Tiny/16 on CIFAR-100: training and validation
  loss and accuracy across 100 epochs.}
  \label{fig:baseline-curves}
\end{figure}

\paragraph{Discussion.}
After epoch 74, the validation accuracy drops slightly instead of
continuing to improve, which suggests mild overfitting in the later
stage of training. This makes the baseline a reasonable starting point:
it is strong enough to be competitive, but it still leaves room for
later methods to improve generalization.
